%%% REVISION NOTES (v2-solo: Shrager sole author, Moews's LT1 work as
%%% postscript) %%%
%%%
%%% Per Alex Magoun's editorial guidance:
%%%   - Cut the IPL-Lisp comparison (old Sec. 2.2 second half and the
%%%     "Reflections: IPL, Lisp, and the Arc of AI" section).
%%%   - Reduced the LLM-debugging material to one short subsection;
%%%     the full narrative is reserved for a possible later submission.
%%%     The Simon's J's / J74 episode is retained because it bears on
%%%     the executable-archaeology argument, not the LLM story.
%%%   - Added Sec. "What Executable Archaeology Contributes" to answer
%%%     Alex's question directly.
%%%   - Moews's IPL-I (LT1) reanimation appears as a Postscript
%%%     section narrating how that work came to light, rather than as
%%%     co-authored content. TODO: confirm with David that he is
%%%     comfortable with this treatment, review the description of his
%%%     work for accuracy, and decide whether he prefers
%%%     acknowledgment vs. co-authorship.
%%%
%%% TODO: The date of the first machine proof. The previous draft said
%%% machine proofs came from IPL-II runs "late 1956 and early 1957";
%%% standard accounts (Simon's Models of My Life ch. 13; Newell & Shaw
%%% 1957) put the first JOHNNIAC proof in August 1956. Text below now
%%% says August 1956 -- verify against the archival record.
%%%
%%% Carried over from v1 notes, still open:
%%%   - Dedicate to Richard Wexelblat?
%%%   - Mention John Hessler's project at Johns Hopkins (prelim LT work)
%%%   - Dick's wonderful paper (at beginning of history)

\documentclass[journal]{IEEEtran}

\usepackage{cite}
\usepackage{amsmath}
\usepackage{graphicx}
\usepackage{url}
\usepackage{hyperref}
\usepackage{booktabs}
\usepackage{listings}

\lstset{
  basicstyle=\ttfamily\small,
  breaklines=true,
  frame=single,
  xleftmargin=1em,
  framexleftmargin=0.5em
}

\begin{document}

\title{Executable Archaeology: Reanimating the Logic Theorist from its
  IPL-V Source}

\author{Jeff~Shrager,~Ph.D.\\
Bennu Climate, Inc.\\
\texttt{jshrager@gmail.com}}

\maketitle

\begin{abstract}

The Logic Theorist (LT), created by Allen Newell, J.~C.~Shaw, and
Herbert Simon in 1955--1956, is widely regarded as the first
artificial intelligence program. While the original conceptual model
was described in 1956, it underwent several iterations as the
underlying Information Processing Language (IPL) evolved. Here I
describe the construction of a new IPL-V interpreter, written in
Common Lisp, and the faithful reanimation of the Logic Theorist from
code transcribed directly from Stefferud's 1963 RAND technical report.
Stefferud's version represents a pedagogical re-coding of the original
heuristic logic into the standardized IPL-V. The reanimated LT
successfully proves 16 of 23 attempted theorems from Chapter~2 of
\textit{Principia Mathematica}, results that are historically
consistent with the original system's behavior within its search
limits. To the author's knowledge, this is the first successful
execution of the original Logic Theorist code in over half a
century. A postscript describes David Moews's complementary
reanimation of the never-machine-executed 1956 IPL-I version, which
came to light after this work was first reported, and which sharpens
the question of what executable archaeology contributes to the history
of AI.

\end{abstract}


\begin{IEEEkeywords}
Logic Theorist, IPL-V, IPL-I, artificial intelligence history,
software archaeology, A.~Newell, H.~A.~Simon, J.~C.~Shaw, list
processing, digital preservation
\end{IEEEkeywords}

\section{Introduction}
\label{sec:intro}

In the summer of 1956, a small group of researchers gathered at
Dartmouth College for a workshop on what John McCarthy had dubbed
``artificial intelligence.'' Among them were Allen Newell and Herbert
Simon, who arrived with something none of the other attendees had: a
program---albeit one not yet running on any machine---that could
discover proofs of mathematical theorems.\footnote{It is worth noting
that Art Samuel's checkers-playing program existed a couple of years
before LT, and it too was heuristic. But Samuel's goal was to build a
checkers player, whereas Newell, Simon, and Shaw's explicit goal was
to model human cognition across many domains. Samuel's effort was
brilliant in its own right---he also built arguably the first machine
learning program, and the first system to learn from self-play.}

The Logic Theorist (LT)\cite{LT1956a,LT1956b} would ultimately prove
38 of the first 52 theorems in Chapter~2 of Whitehead and Russell's
\textit{Principia Mathematica}~\cite{whitehead1910principia}, in some
cases finding proofs more elegant than those produced by the human
authors. For Theorem~2.85, LT discovered a shorter, more direct proof
than the one published in \textit{Principia}. Simon showed the new
proof to Bertrand Russell himself, who ``responded with
delight''~\cite[p.~167]{mccorduck2004machines}. Edward Feigenbaum,
then an undergraduate taking a graduate course from Simon, recalls the
moment the project's significance was first announced. As he later
told Pamela McCorduck: ``It was just after Christmas vacation---January
1956---when Herb Simon came into the classroom and said, `Over
Christmas Allen Newell and I invented a thinking machine.' And we all
looked blank''~\cite[p.~138]{mccorduck2004machines}.

Beyond reproducing \textit{Principia} proofs, the Logic Theorist
introduced heuristic search strategies---means--ends analysis and
subgoal decomposition---derived from protocol analyses of human
problem solving. These ideas later became foundational in both
artificial intelligence and cognitive architectures such as the
General Problem Solver (GPS)~\cite{newell1959gps}, Soar, ACT, and
EPAM~\cite{feigenbaum1961epam}. LT thus sits at the very beginning of
the fusion of AI and cognitive psychology.

LT was implemented in the Information Processing Language (IPL), a
language created by the same team---Newell, Shaw, and
Simon---specifically for building cognitive
models~\cite{newell1964ipl}. IPL introduced list manipulation, dynamic
memory allocation, symbols as first-class objects, recursion, and
higher-order functions, and was an explicit, acknowledged ancestor of
Lisp~\cite{mccarthy1960lisp} and of other early list-processing
systems such as Weizenbaum's SLIP. Yet while Lisp has thrived for
nearly seven decades, IPL has been extinct since the mid-1960s.
Although numerous IPL-V implementations existed in the 1960s on
various machines, no maintained implementations appear to be available
in modern environments. The only prior modern attempt known to the
author is Richman's unpublished IPL-V interpreter, written in
Microsoft BASIC, which he used to run EPAM~IV in the early
1990s~\cite{richman1995epam}.\footnote{Richman wrote to the author
(personal communication, Sep~26, 2025, quoted with permission): ``I
just interpreted IPL-V in Microsoft Basic using the IPL-V manual that
I found at the CMU library.''}

This paper describes the construction of a complete IPL-V interpreter
in Common Lisp, and the faithful reanimation of the Logic Theorist
from code transcribed directly from Stefferud's 1963 RAND technical
report~\cite{stefferud1963lt}. More than 99\% of the LT code running
in this reconstruction is the original code from that report. The
reanimated LT proves 16 of 23 attempted theorems from
\textit{Principia Mathematica}, results that are historically
consistent with the original LT's known behavior.

This work illustrates a methodological approach that might be called
\emph{executable archaeology}: the study of historically significant
computational systems through faithful reconstruction and execution of
their original code. Early AIs were not merely theoretical proposals
but functioning programs whose detailed behavior cannot be fully
understood from published descriptions alone. Reanimating such systems
provides a new empirical window into the origins of artificial
intelligence; Section~\ref{sec:contribution} takes up exactly what
that window shows, and the Postscript
(Section~\ref{sec:postscript}) describes a complementary reanimation,
by David Moews, of the program's never-machine-executed 1956
ancestor, which came to light after this work was first reported.

In prior work, a team including the present author restored
Weizenbaum's original ELIZA to operation on a reconstructed CTSS
running on an emulated IBM 7094~\cite{lane2025eliza}. Although the
Logic Theorist and ELIZA addressed very different problems---formal
theorem proving versus natural-language conversation---both emerged
from the same symbolic, list-processing programming lineage
originating in IPL. However, whereas the ELIZA work rested on a stack
that emulated vintage hardware and operating systems, the LT
reanimation emulates the abstract IPL-V machine.

\section{Historical Background}
\label{sec:history}

\subsection{The Origins and Versions of the Logic Theorist}

The story of the Logic Theorist begins in the early 1950s at the RAND
Corporation. Herbert Simon and Allen Newell, collaborating with
J.~C. (Cliff) Shaw, sought to simulate complex thought by manipulating
symbols. However, the history of the program is characterized by a
series of shifting specifications across different versions of the
Information Processing Language.

The system described in the famous 1956 RAND report
(P-868)~\cite{LT1956a}, and in nearly identical form in the IRE
Transactions that September~\cite{LT1956b}, was a ``paper'' version,
written in what was called IPL-I or the ``Logic Language'' (LL).
IPL-I was never implemented on any machine: it was a pseudocode for an
imagined abstract machine, executed entirely by hand. At the time of
the Dartmouth workshop in the summer of 1956, Newell and Simon noted
that the program was ``essentially specified'' and that hand
simulations suggested it would prove the 60-odd theorems of Chapter~2
of \textit{Principia Mathematica}. These hand simulations famously
employed ``the human computer'': Simon's wife, his children, and
graduate students acting as program
components~\cite[Ch.~13]{simon1996models}.

The first version to run on real hardware was written in IPL-II for
the RAND JOHNNIAC~\cite{gruenberger1968johnniac,newellshaw1957},
producing its first machine proof in August
1956~\cite{simon1996models,newellshaw1957}. The widely cited result
that LT proved 38 of the first 52 theorems in
\textit{Principia}---including the discovery of a more elegant proof
for Theorem 2.85 than the one produced by Whitehead and
Russell---stems from the empirical explorations reported in
1957~\cite{LT1957}.

The code used in this reanimation, transcribed from Stefferud's 1963
memorandum, represents a third distinct phase. The program was
converted to IPL-V by Fred Tonge and further developed by Stefferud as
a pedagogical re-coding of the original heuristic logic into IPL-V,
the first version of the language made widely available for public
use~\cite{stefferud1963lt}. Thus, while this reanimation is a faithful
execution of original code published by the RAND team, it is a
restoration of the 1963 ``modernized'' version of the 1956 conceptual
model. (The Postscript describes the subsequent reanimation of the
1956 version itself.)

%%% TODO: confirm/adjust years in the 2026 rows.
\begin{table}[h]
\centering
\caption{Evolution of the Logic Theorist and IPL}
\label{tab:lt_versions}
\begin{tabular}{@{}llll@{}}
\toprule
\textbf{Year} & \textbf{Version} & \textbf{Platform} & \textbf{Historical Milestones} \\ \midrule
1956 & IPL-I (LL) & Manual & Conceptual spec; hand-simulated only. \\
1956--57 & IPL-II & JOHNNIAC & First machine proofs; 38 of 52 theorems. \\
1958 & IPL-IV & IBM 704 & Internal RAND use; transition to IPL-V. \\
1963 & IPL-V & Various & Stefferud's pedagogical LT implementation. \\
2026 & IPL-V & Lisp & This reanimation (Stefferud's source). \\
2026 & IPL-I & Python & Moews: first machine execution (Postscript). \\ \bottomrule
\end{tabular}
\end{table}

\subsection{IPL: The Invention of Symbol- and List-Processing}

To implement LT, Newell, Shaw, and Simon needed a programming language
suited to symbolic reasoning. No such language existed, so they
created the ``Information Processing Language.'' As noted, IPL went
through several iterations. The earliest versions (IPL-I and IPL-II)
powered the 1956 hand simulations and subsequent JOHNNIAC runs;
IPL-V, released in the early 1960s~\cite{newell1964ipl}, became the
definitive public standard.

IPL-V runs on an abstract machine organized around \textit{cells}.
Each cell has an associated symbol. There is no single distinguished
stack; every cell is potentially a stack of symbols that can be pushed
and popped. Symbols are central to IPL-V: they are first-class
objects, and every cell, data structure, routine, and control element
is addressable by a symbol (its name). A set of control cells
(H0--H5) governs execution: H0 serves as the parameter and result
``communication'' cell, H1 serves as both program counter and call
stack, H2 is the free list for dynamic memory allocation, and H3 is
the cycle counter. Ten working cells (W0--W9) provide scratch
storage. Everything in IPL-V---data, routines, even the program
itself---is a list of cells, and thus also a list of symbols. Routines
are invoked by pushing their symbol (name) onto H1. The IPL machine
interprets the operation indicated in that cell, and then continues to
whatever is the indicated next cell. Elegantly, because H1 is both
program counter and call stack, no special machinery is required to
operate the machine aside from basic stack manipulation instructions.

The language's built-in operations are called J-functions (because
their names begin with ``J'': J0, J1, J2, and so on). There are
approximately 150 of these, implementing operations from basic list
manipulation to ``generator'' (iterator) protocols. In the original
IPL-V, many J-functions were themselves written in IPL, although some
were implemented at the level of the underlying physical machine.

\section{Source Materials and Transcription}
\label{sec:sources}

Unlike the ELIZA reanimation, which began with the dramatic discovery
of original source code printouts in Weizenbaum's archives at
MIT~\cite{lane2025eliza}, the source materials for the Logic Theorist
were already publicly available. Two documents were indispensable:

\begin{enumerate}
\item \textbf{Stefferud (1963)}, \textit{The Logic Theory Machine: A
  Model Heuristic Program}, RAND Corporation memorandum
  RM-3731~\cite{stefferud1963lt}. This report contains a complete
  listing of the LT program in IPL-V---nearly 3,000 lines---along with
  detailed descriptions of every aspect of LT's workings, including
  numerous flow charts.

\item \textbf{Newell et al.\ (1964)}, \textit{Information Processing
  Language-V Manual}, 2nd edition~\cite{newell1964ipl}. The definitive
  reference for IPL-V, including the semantics of all J-functions,
  cell structure conventions, and the generator protocol, as well as
  numerous student exercises and sample programs.
\end{enumerate}

Both documents proved to be exceptionally well-written. There were,
however, occasional ambiguities---places where the specification
assumed familiarity with conventions of 1960s computing (such as BCD
character encoding) that are no longer common knowledge.

The LT code was transcribed from Stefferud's report into a Google
spreadsheet by the author and Anthony Hay. Subsequently, Rupert Lane
identified several transcription errors using his ``gridlock''
tool~\cite{lane2024gridlock}, software designed to accurately extract
tabular code from PDF documents. Despite multiple people putting
multiple sets of eyes on the code, a surprising number of small
transcription errors persisted through numerous rounds of
review.\footnote{I probably shouldn't have been surprised by this
given the number of typos that have persisted through numerous rounds
of review of this paper!}

The transcribed code was converted into an S-expression format
(\texttt{.liplv} files) suitable for loading by the Lisp-based
interpreter. Conveniently, Lisp comment characters work in these
files, allowing corrections and annotations to be documented inline.

\section{A New IPL-V Interpreter in Lisp}
\label{sec:interpreter}

I implemented a new IPL-V interpreter to execute the Logic Theorist
listing. The interpreter implements the symbolic data structures and
execution model described in the IPL-V manual and supports the
primitives required by the program. Although the earliest versions of
the Logic Theorist ran on the RAND JOHNNIAC, by the time the Stefferud
listing, used here, was published (1963) IPL-V had already been
implemented on multiple machines. Accordingly, this reconstruction
executes the program within a machine-independent IPL-V interpreter
rather than attempting to reproduce a specific JOHNNIAC hardware
environment. The full interpreter and transcription of the Logic
Theorist code are available in the repository associated with this
paper.

I chose to emulate the IPL-V abstract machine, rather than emulating
the original hardware environment (as was done for ELIZA on the IBM
7094 and CTSS~\cite{lane2025eliza}), for several reasons. First, as
mentioned just above, the code I am using is IPL-V, and so is not
JOHNNIAC-specific. Additionally, no thorough JOHNNIAC emulator exists,
and in any case we do not have the JOHNNIAC IPL code. There does exist
an IBM 7094 emulator that includes an IPL-V interpreter running under
a variant of the IBSYS batch operating system \cite{ibsys}, and this
emulator includes demo scripts for IPL-V. However, this IPL-V
interpreter is a binary with no source code and no documentation. I
attempted to run LT on this platform but it failed early in execution
with no useful information about the cause. Because there is no source
code, this failure was essentially undebuggable.

More fundamentally, however, one of my central goals in undertaking
this project was not merely to get LT to run, but to
\textit{elucidate} IPL-V---to develop a deep understanding of the
language, the machine model, and the infrastructure of the
period. Writing my own interpreter was essential to that goal. IPL-V
is defined as an abstract machine specification, not as a program for
a specific computer; it was implemented on many different machines
through the mid-1960s before being supplanted by Lisp.

My interpreter (\texttt{iplv.lisp}) implements the complete IPL-V
abstract machine as described in the 1964 manual: the cell table, the
control cells H0--H5, the working cells W0--W9, and the push-down
mechanisms. Lists are stored as linked chains of cells, following the
manual's specification. It is worth noting that the push-down (stack)
mechanism is itself implemented as linked lists within the IPL-V
system. Push and pop operations, which appear to be hardware
primitives of the abstract machine, are internally built on the same
list mechanism that underlies everything else in IPL-V, including data
structures, routines, and IPL programs themselves; like Lisp, IPL-V is
lists all the way down.

The core of my interpreter is a function, \texttt{ipl-eval}, that
fetches and executes instructions by executing the list named in the
program counter (H1). My IPL evaluator is re-entrant---a capability
trivially implementable in Lisp---supporting the recursive calls that
IPL-V requires. The J-functions are implemented as Lisp
functions. While the interpreter implements the full IPL-V machine,
only the subset of J-functions required by LT (and a few test
programs) is currently implemented, amounting to several dozen of the
approximately 150 defined in the manual.

A complete run of LT attempting proofs of 23 theorems is available at
\cite{runoutput}. This takes approximately 574,000 IPL machine cycles
and creates over 576,000 symbols\footnote{Yeah, I was surprised
too. It's really just a coincidence!}. It finishes in a few seconds on
a typical modern laptop. On the JOHNNIAC at RAND, in the mid 1950s,
this likely would have required hours.

\section{Debugging Against the Archival Record}
\label{sec:debugging}

%%% Note: the extended account of LLM-assisted debugging has been cut
%%% per editorial guidance and may seed a separate submission. The
%%% Simon's J's episode is retained because it bears directly on the
%%% executable-archaeology argument.

The most difficult phase of debugging began once the interpreter was
stable enough that LT ran without crashing: the program would prove
trivial theorems but fail on simple ones, fall into runaway
recursions, or fail to recognize a completed proof---anomalies whose
causes were far from their symptoms. This experience had an uncanny
resonance with the original developers' accounts. In an interview
conducted by Pamela McCorduck for her pioneering history of AI, Cliff
Shaw described the most daunting phase of troubleshooting as occurring
late in development, when the system was beyond crashes and produced
behavior that appeared superficially reasonable, leaving the
developers in a quandary as to whether the results were valid
consequences of the heuristic logic or actual bugs, in a memory that
had become a churning and impenetrable web of interconnected
lists~\cite{cmuarchiveshawinterview}. Sixty years later, with the same
program, I faced the identical challenge. In this latter phase I used
large language models (Google's Gemini and Anthropic's Claude) as
debugging assistants, performing the grinding work of reading
thousands of lines of execution traces and tracking machine state
across tens of thousands of cycles; a fuller account of that
collaboration is beyond the scope of this paper.

One episode, however, bears directly on the methodology at issue
here. The final substantive bug was localized to J74, the list-copying
function, yet my Lisp implementation appeared to match the manual's
prose description; the specification itself seemed to support the
buggy behavior. Resolution came not from documentation but from an
artifact: ``Simon's J's,'' a set of original IPL-V punched cards from
1962 preserved by the Computer History
Museum~\cite{bochannek2012simonsjs}, representing the production
implementation of the J-functions at the moment IPL-V was being
standardized. Comparing the 1962 card-deck implementation of J74 with
mine confirmed a hypothesized discrepancy, and a one-line patch made
LT run completely correctly. The manual's prose, my careful reading of
it, and inference from my code had all been insufficient; the original
1962 source was the ground truth.

\section{What the Logic Theorist Proves}
\label{sec:results}

The Logic Theorist attempts to prove theorems in propositional logic
using the five axioms of \textit{Principia Mathematica}, Chapter~2.
The notation used by the original authors is given in
Table~\ref{tab:notation}.

\begin{table}[h]
\centering
\caption{Logical notation used by the Logic Theorist}
\label{tab:notation}
\begin{tabular}{cl}
\toprule
\textbf{Symbol} & \textbf{Meaning} \\
\midrule
\texttt{I} & Implication ($\rightarrow$) \\
\texttt{V} & Disjunction ($\lor$) \\
\texttt{*} & Conjunction ($\land$) \\
\texttt{-} & Negation ($\lnot$) \\
\texttt{=} & Biconditional ($\leftrightarrow$) \\
\texttt{.=.} & Definitional equality \\
\bottomrule
\end{tabular}
\end{table}

LT employs several proof methods:

\begin{itemize}
\item \textbf{Substitution}---substitute variables in an axiom or
  previously proved theorem to match the goal.
\item \textbf{Detachment} (modus ponens)---from $A \rightarrow B$ and
  $A$, conclude $B$.
\item \textbf{Forward chaining}---apply substitution and detachment
  working forward from known theorems.
\item \textbf{Backward chaining}---work backward from the goal,
  seeking premises that would yield it.
\item \textbf{Sublevel replacement}---replace a subexpression using a
  known equivalence.
\end{itemize}

Each proof attempt is bounded by a settable effort limit in terms of
IPL machine cycles. (In all below examples this limit is set to 20,000
cycles.) This is a soft cap: it prevents the system from
\textit{beginning} a new subproblem exploration once exceeded, but
does not interrupt a search in progress.

Table~\ref{tab:axioms} shows the axioms and definitions that are
initially provided to LT for the runs described in the present
paper. Importantly, as LT proves a given theorem, that theorem is
treated as an axiom and can be used to support subsequent
proofs. Thus, LT can be said, in a sense, to learn!

\begin{table}[h]
\centering
\caption{Axioms and definitions}
\label{tab:axioms}
\begin{tabular}{ll}
\toprule
\textbf{Name} & \textbf{Axiom/Def}\\
\midrule
1.01  & \texttt{((PIQ).=.(-PVQ))}\\
2.33  & \texttt{((PVQVR).=.((PVQ)VR))}\\
3.01  & \texttt{((P*Q).=.-(-(PV-Q))}\\
4.01  & \texttt{((P=Q).=.((PIQ)*(QIP)))}\\
1.2   & \texttt{((AVA)IA)}\\
1.3   & \texttt{(BI(AVB))}\\
1.4   & \texttt{((AVB)I(BVA))}\\
1.5   & \texttt{((AV(BVC))I(BV(AVC)))}\\
1.6   & \texttt{((BIC)I((AVB)I(AVC)))}\\
\bottomrule
\end{tabular}
\end{table}

Table~\ref{tab:results} shows the results of my reanimated LT on the
23 theorems attempted: 16 are proved, and the remaining 7 are
not---results that are historically consistent with the original LT's
known behavior, given an effort limit of 20,000 IPL machine cycles.

\begin{table}[h]
\centering
\caption{Theorems attempted}
\label{tab:results}
\begin{tabular}{lll}
\toprule
\textbf{Name} & \textbf{Formula} & \textbf{Result} \\
\midrule
2.01 & \texttt{(PI-P)I-P} & Proved \\
2.02 & \texttt{QI(PIQ)} & Proved \\
2.04 & \texttt{(PI(QIR))I(QI(PIR))} & Proved \\
2.05 & \texttt{(QIR)I((PIQ)I(PIR))} & Proved \\
2.06 & \texttt{(PIQ)I((QIR)I(PIR))} & Proved \\
2.07 & \texttt{PI(PVP)} & Proved \\
2.08 & \texttt{PIP} & Proved \\
2.10 & \texttt{-PVP} & Proved \\
2.11 & \texttt{PV-P} & Proved \\
2.12 & \texttt{PI--P} & Proved \\
2.13 & \texttt{PV---P} & Proved \\
2.14 & \texttt{--PIP} & Not proved \\
2.15 & \texttt{(-PIQ)I(-QIP)} & Not proved \\
2.20 & \texttt{PI(PVQ)} & Proved \\
2.21 & \texttt{-PI(PIQ)} & Not proved \\
2.24 & \texttt{PI(-PVQ)} & Proved \\
3.13 & \texttt{(-(P*Q))I(-PV-Q)} & Not proved \\
3.14 & \texttt{(-PV-Q)I(-(P*Q))} & Not proved \\
3.24 & \texttt{-(P*-P)} & Not proved \\
4.13 & \texttt{P=--P} & Not proved \\
4.20 & \texttt{P=P} & Proved \\
4.24 & \texttt{P=(P*P)} & Proved \\
4.25 & \texttt{P=(PVP)} & Proved \\
\bottomrule
\end{tabular}
\end{table}

A representative proof trace is shown below. Here LT proves Theorem
2.01, \texttt{(PI-P)I-P}:

\begin{lstlisting}
2.01    (PI-P)I-P

PROOF FOUND.

    GIVEN           *1.2  (AVA)IA
    SUBSTITUTION    .0    (PI-P)IP
    SUBLEVEL REPL   2.01  (PI-P)I-P
    Q.E.D.

EFFORT      LIMIT 20000   ACTUAL 5579
SUBPROBLEMS LIMIT 50      ACTUAL 1
SUBSTITUTIONS LIMIT 50    ACTUAL 2
\end{lstlisting}

The proof begins from Axiom~*1.2, \texttt{(AVA)IA}, substitutes
\texttt{-P} for \texttt{A} to obtain \texttt{(-PV-P)I-P}, then applies
sublevel replacement (using the fact that \texttt{P} implies
\texttt{PVP}) to transform the antecedent, yielding the desired
theorem. The Appendix depicts a more complex proof (4.25), and the
original output, including both of these proofs, is at
\cite{runoutput}.

\section{What Executable Archaeology Contributes}
\label{sec:contribution}

%%% This section is new, written in response to Alex's question:
%%% "What exactly does it contribute to understanding the origins of
%%% AI in this case?" It draws on both the present LT5 work and
%%% Moews's LT1 work described in the Postscript.

It is reasonable to ask what is learned by running this old code that
could not be learned by reading it, or by reading the abundant
secondary literature on the Logic Theorist. The present case---taken
together with the complementary reanimation described in the
Postscript---suggests several answers.

\textbf{It converts historical claims into checkable ones.} The 1956
report claimed, on the basis of hand simulation, that the program
would prove the theorems of \textit{Principia} Chapter~2; the 1957
paper reported the famous 38-of-52 result for the JOHNNIAC
version. Until now, such claims rested on testimony and on archival
output listings. Reanimation allows them to be tested directly: the
reconstructed LT proves 16 of 23 attempted theorems under a fixed
effort limit, a result consistent with the original
reports. Moreover, the reconstruction is now an experimental
instrument: effort limits, axiom orderings, and theorem sequences can
be varied, and the consequences observed on the genuine historical
artifact rather than on a modern paraphrase of it.

\textbf{It shows that some semantics survive only in artifacts, not
in documentation.} The J74 episode (Section~\ref{sec:debugging}) is a
small but pointed demonstration: the IPL-V manual's prose description
of a core primitive was consistent with an incorrect implementation,
and the discrepancy could be resolved only by consulting the original
1962 punched-card source preserved by the Computer History Museum.
For the history of computing this is a methodological moral:
specifications and manuals, however well written, systematically
under-determine the systems they describe, and executable
reconstruction is one of the few ways to discover where.

\textbf{It corrects and annotates the published record.} Faithful
reconstruction imposes a discipline that reading does not: every typo,
ambiguity, and inconsistency in the source materials must be found,
repaired, justified, and recorded, because the machine will not
silently forgive them. The transcription hazards encountered in
working from the 1963 report---and, as the Postscript describes, the
outright errors in the 1956 printed listing---are now documented in
the project repositories. These repairs are the ``archaeology'' in
executable archaeology.

\textbf{It makes the gap between cognitive theory and implementation
measurable.} This contribution emerges most clearly from the
juxtaposition of the present work with Moews's reanimation of the
1956 IPL-I version (Postscript). The 1956 program---roughly 400
lines---is something close to a direct expression of Newell and
Simon's hypothesized cognitive operators for human theorem proving: a
high-level specification of a theory of cognition. It could afford to
be that, because its ``hardware'' was human---intelligent, flexible,
and able to silently fill in whatever the specification left
unsaid. The Stefferud listing---nearly 3,000 lines---is what the same
design became when Shaw and his successors had to make a real,
literal-minded machine do the whole thing end-to-end. Executing both
versions converts a familiar qualitative observation about early
AI---that getting from idea to running program was most of the
work---into something that can be measured, localized, and studied
line by line. Most strikingly, as Moews discovered, the 1956 program
contains no machinery at all to assemble or print the proofs it
finds: that labor was silently delegated to the human simulators. The
omission is itself historical evidence---it marks precisely a
boundary of what Newell and Simon considered the \emph{theory}
(finding the proof) versus mere clerical work (writing it down).

\section{Related Work}
\label{sec:related}

The reanimation of historical software systems has gained momentum in
recent years. Most directly related is the ELIZA reanimation by Lane
et al.~\cite{lane2025eliza}, in which a team including the present
author restored Weizenbaum's original ELIZA to operation on a
reconstructed CTSS running on an emulated IBM~7094. That project
reconstructed the complete original hardware and operating system
stack; the present work instead emulates the abstract machine
specification in a modern language.

Reimplementing the Logic Theorist's \textit{algorithms} in a modern
language is a relatively straightforward exercise---essentially an
introductory AI course project, and one that has surely been done many
times. A 1968 paper, for example, describes an implementation of LT in
LISP~1.5~\cite{millstein1968logic}. The present work is a
fundamentally different undertaking: rather than rewriting LT's logic
in a modern language, it constructs a faithful emulator of the IPL-V
abstract machine and runs the original code, transcribed directly from
a 1963 technical report. The distinction is significant: a Lisp
reimplementation demonstrates that LT's algorithms can be expressed in
Lisp (which is unsurprising), while the present work demonstrates that
the original IPL-V code, as written by Newell, Simon, and Shaw's team,
functions as described, and, moreover, enables us to perform
experiments on the real, working system of machine and program.

\section{Next Steps}
\label{sec:next}

The IPL-V interpreter developed for this project is not specific to
the Logic Theorist. It implements the general IPL-V abstract machine
and can, in principle, run any IPL-V program. I am in the process of
building a corpus of early IPL-V AI programs. The most significant
targets are the General Problem Solver (GPS)~\cite{newell1959gps} and
EPAM (Elementary Perceiver and
Memorizer)~\cite{feigenbaum1961epam}. EPAM source code is available in
the Feigenbaum archives at Stanford, and archival work to locate
complete source code for GPS is in progress. Now that a proven IPL-V
interpreter exists, running these programs should require
substantially less effort than the initial construction of the
interpreter itself.

The interpreter, the transcribed LT code, and complete outputs as
described in this paper are open sourced at
\url{https://github.com/jeffshrager/IPL-V}.

\section{Postscript: The 1956 Program Runs At Last}
\label{sec:postscript}

%%% TODO: David has reviewed/should review this section for accuracy.
%%% Facts below are drawn from his email of May 10, 2026 and his
%%% repository README; placeholders marked TODO (DM) need his input.

Shortly after a preprint of this paper appeared on arXiv, I received
an email from David Moews, a mathematician previously unknown to
me. He had, it turned out, just been working on the same
problem---but had attacked it from the other end of the program's
history. Rather than reimplementing IPL-V, Moews had built an
interpreter, in Python, for the original IPL-I pseudocode itself, as
published in the 1956 RAND report P-868~\cite{LT1956a} and, in nearly
identical form, in Newell and Simon's 1956 IRE
paper~\cite{LT1956b}. He had the code working in February
2026~\cite{moewsrepo}.

The significance of this bears emphasis. As discussed in
Section~\ref{sec:history}, IPL-I was never implemented on any machine:
it was hand-executed by Simon's family and students simulating the
imagined IPL-I machine. Every prior execution of this program in
history was performed by human beings. Moews's run is therefore, to
the best of our joint knowledge, the first time the IPL-I version of
the Logic Theorist has ever been executed by a machine---some seventy
years after it was written, and pushing the window of ``RetroAI''
reanimation back to the very first AI program ever specified.

Moews reports that the printed listing contains typos and other
problems and is not runnable exactly as published; it required minor
patches and bug fixes, each documented in his
repository~\cite{moewsrepo}. One omission he found is particularly
telling. Although the 1956 code is capable of finding individual proof
steps, and of determining when enough steps have been found to
constitute a proof, it contains no logic whatsoever to print the steps
or to assemble them into the actual proof. The human simulators, of
course, needed no such logic---they could see the steps they had
executed and write the proof down. Moews's solution was to have the
interpreter print each potential proof step as it is found, and to
write a separate program that assembles and checks the proof from the
resulting set of steps. He also built a driver that imitates the 1957
``Empirical Explorations'' regime~\cite{LT1957}, running the program
successively on the theorems of \textit{Principia} Chapter~2 and
allowing each proof to use all axioms and all previous Chapter~2
theorems.

%%% TODO (DM): a sentence or two on results -- which theorems the
%%% IPL-I version proves, and how that compares to Table 4 and to the
%%% 1956/1957 reports.

The two reconstructions bracket the Logic Theorist's history: the 1956
program is the conceptual model as first specified, in roughly 400
lines that hew closely to Newell and Simon's hypothesized cognitive
operators; the 1963 program reanimated in the body of this paper is
what that model became---nearly 3,000 lines---once it had to run,
end-to-end, on a literal-minded machine. That the two efforts were
undertaken independently, by researchers unknown to each other, and
converged on complementary working reconstructions of adjacent stages
of the same program, lends each a measure of replication: agreement
between them is evidence that the reconstructions reflect the
historical system rather than either reconstructor's assumptions. The
implications of this juxtaposition for what executable archaeology can
teach us are taken up in Section~\ref{sec:contribution}; Moews's
interpreter, repaired source, and proof-assembly tools are available
at \url{https://github.com/dmoews/logic-theorist}.

\section*{A Personal Coda}

I studied under both Newell and Simon at Carnegie Mellon University in
the 1980s, where my work focused on modeling the cognition of
scientists and engineers. IPL and LT were long since history by that
time; I unfortunately never discussed either with either Simon or
Newell---we barely learned about these pioneering efforts in our
classes. But the intellectual values they instilled---rigor about
mechanism, respect for the complexity of cognition, and the conviction
that thinking could be understood computationally---shaped the rest of
my career. Now, at the end of that career, I find myself studying
Newell, Simon, and Shaw \textit{as} scientist-engineers:
reverse-engineering their design decisions, reconstructing their
reasoning, and trying to understand how they invented AI and cognitive
science. The subject of my training---the cognition of scientists and
engineers---and its object have folded into one another.

\section*{Acknowledgements}

The following people (listed in no particular order) contributed to
this project in ways ranging from code transcription to weekly
encouragement: Anthony Hay, Arthur Schwarz, Leigh Klotz, David
M.\ Berry, Rupert Lane, Amy Majczyk, Emily Davis, Ethan Ableman, and
Paul McJones. I am particularly grateful to Hay, Schwarz, Klotz, and
Berry, who put up with weekly progress updates throughout the project
and responded with encouragement that was always helpful and sometimes
technically substantive. Rupert Lane's gridlock software caught
several transcription errors that had survived multiple rounds of
human review.

I am especially grateful to David Moews for generously sharing his
IPL-I reanimation, for reviewing the Postscript describing it, and for
the correspondence that has enriched this paper. I also thank the
Computer History Museum for preserving and making available ``Simon's
J's,'' the original 1962 IPL-V J-function card deck, which proved
essential during debugging, and Alexander Magoun and David Walden for
editorial guidance.

Finally, I want to acknowledge Allen Newell and Herbert Simon, under
whom I studied at CMU in the 1980s. The subject of that
training---the cognition of scientists and engineers---has, in this
project, become inseparable from its object.

{\footnotesize Portions of the debugging effort, and portions of the
drafting of this paper, were carried out with the aid of various LLMs,
but all concepts and final editorial decisions were the author's
alone.}

\begin{thebibliography}{99}

\bibitem{LT1956a} A. Newell and H. A. Simon, \textit{The Logic Theory
  Machine -- A Complex Information Processing System}, RAND Corp.,
  Santa Monica, CA, USA, Rep. P-868, 1956.

\bibitem{LT1956b} A. Newell and H. A. Simon, ``The logic theory
  machine--A complex information processing system,'' \textit{IRE
  Trans. Inf. Theory}, vol. 2, no. 3, pp. 61--79, Sep. 1956.

\bibitem{LT1957} A. Newell, J. C. Shaw, and H. A. Simon, ``Empirical
  explorations with the logic theory machine: A case study in
  heuristics,'' in \textit{Proc. Western Joint Comput. Conf.}, 1957,
  pp. 218--230.

\bibitem{newellshaw1957} A. Newell and J. C. Shaw, ``Programming the
  logic theory machine,'' in \textit{Proc. Western Joint
  Comput. Conf.}, 1957, pp. 230--240.

\bibitem{whitehead1910principia} A. N. Whitehead and B. Russell,
  \textit{Principia Mathematica}. Cambridge, UK: Cambridge
  Univ. Press, 1910--1913.

\bibitem{newell1964ipl} A. Newell, et al. \textit{Information
  Processing Language-V Manual}, 2nd ed. Englewood Cliffs, NJ, USA:
  Prentice-Hall, 1964.

\bibitem{mccarthy1960lisp} J. McCarthy, ``Recursive functions of
  symbolic expressions and their computation by machine, Part I,''
  \textit{Commun. ACM}, vol. 3, no. 4, pp. 184--195, Apr. 1960.

\bibitem{simon1996models} H. A. Simon, \textit{Models of My Life}.
  Cambridge, MA, USA: MIT Press, 1996.

\bibitem{cmuarchiveshawinterview}
  Pamela McCorduck Collection, Carnegie Mellon University Archives.

\bibitem{ibsys} R. Storey, ``B7094'' [Online]. Available:
  \url{https://github.com/Bertoid1311/B7094} (Accessed: Mar. 11,
  2026).

\bibitem{stefferud1963lt} E. Stefferud, ``The Logic Theory Machine: A
  Model Heuristic Program,'' RAND Corp., Santa Monica, CA, USA,
  Rep. RM-3731-CC, 1963.

\bibitem{lane2025eliza} R. Lane, A. Hay, A. Schwarz, D. M. Berry, and
  J. Shrager, ``ELIZA Reanimated: Restoring the mother of all chatbots
  to one of the world's first time-sharing systems,'' \textit{IEEE
    Ann. Hist. Comput.}, vol. 47, no. 1, pp. 8--23, 2025.

\bibitem{gruenberger1968johnniac} F. Gruenberger, ``The History of the
  JOHNNIAC,'' RAND Corp., Santa Monica, CA, USA, Rep. RM-5654-PR,
  1968.

\bibitem{mccorduck2004machines} P. McCorduck, \textit{Machines Who
  Think}, 2nd ed. Natick, MA, USA: A. K. Peters, 2004.

\bibitem{bochannek2012simonsjs} A. Bochannek, ``Simon's J's,''
  Computer History Museum Blog, Sep. 26, 2012. [Online]. Available:
  \url{https://computerhistory.org/blog/simons-js/} (Accessed:
  Mar. 11, 2026).

\bibitem{lane2024gridlock} R. Lane, ``Gridlock: Accurate code
  extraction from PDFs,'' 2024. [Online]. Available:
  \url{https://github.com/rupertl/gridlock}

\bibitem{millstein1968logic} J. S. Millstein, ``The Logic Theorist in
  LISP,'' \textit{Comput. J.}, vol. 11, no. 1, pp. 39--44, Feb. 1968.

\bibitem{newell1959gps} A. Newell, J. C. Shaw, and H. A. Simon,
  ``Report on a general problem-solving program,'' in
  \textit{Proc. Int. Conf. Inform. Process.}, 1959, pp. 256--264.

\bibitem{feigenbaum1961epam} E. A. Feigenbaum, ``The simulation of
  verbal learning behavior,'' in \textit{Proc. Western Joint
    Comput. Conf.}, 1961, pp. 121--132.

\bibitem{richman1995epam} H. B. Richman, J. J. Staszewski, and
  H. A. Simon, ``Simulation of expert memory using EPAM IV,''
  \textit{Psychol. Rev.}, vol. 102, no. 2, pp. 305--330, 1995.

\bibitem{moewsrepo} D. Moews, ``Recreation of the 1956 IPL-I version
  of the Logic Theorist theorem prover,'' 2026. [Online]. Available:
  \url{https://github.com/dmoews/logic-theorist} (Accessed: Jun. 9,
  2026).

\bibitem{runoutput} J. Shrager, Output from the reanimated logic
  theorist, [Online]. Available:
  \url{https://github.com/jeffshrager/IPL-V/blob/master/major_results/20260309_CompleteAndCorrect.drb}
  (Accessed: Mar. 11, 2026).

\end{thebibliography}

\clearpage % Ensure all previous material is printed
\onecolumn % Switch to single column mode for the appendix

\appendix

\section{Example Complex Logic Theorist Proof}

To illustrate a complex proof generated by the reconstructed Logic
Theorist, the following derivation corresponds to theorem *4.25 of
\textit{Principia Mathematica}:

\[ P = (P \lor P) \]

This states the idempotence of disjunction. In the LT representation,
the equivalence operator is defined using implication and conjunction:

\[ (P = Q) \equiv ((P \rightarrow Q) \land (Q \rightarrow P)) \]

The Logic Theorist proof trace proceeds as follows:

\begin{enumerate}

\item \textbf{Expand the definition of equivalence:}
LT begins by reducing the goal to a conjunction of two implications:
\begin{quote}
\begin{verbatim}
4.25    P=(PVP)
.0   (PI(PVP))*((PVP)IP)  4.01  ,0 REPLACEMENT
\end{verbatim}
\end{quote}
\textit{Logic form:}
\begin{itemize}
    \item[] $P = (P \lor P) \Rightarrow (P \rightarrow (P \lor P)) \land ((P \lor P) \rightarrow P)$ 
\end{itemize}

\item \textbf{Identify the first implication:}
The system identifies a forward-chaining path to establish the first half of the conjunction:
\begin{quote}
\begin{verbatim}
PI(PVP)    4.20  ,0 FORWARD CHAINING
\end{verbatim}
\end{quote}
\textit{Logic form:} 
\begin{itemize}
    \item[] $P \rightarrow (P \lor P)$
\end{itemize}

\item \textbf{Find the proof via substitution:}

LT utilizes Theorem *2.20 from its memory as the basis for
substitution: \textit{Theorem 2.20 is not given as an axiom, but is a
previously proved theorem that it remembered and applied here; See
comment below.}

\begin{quote}
\begin{verbatim}
GIVEN      2.20      AI(AVB)
SUBSTITUTION  .0    PI(PVP)
\end{verbatim}
\end{quote}
\textit{Logic form:} 
\begin{itemize}
    \item[] The general theorem $A \rightarrow (A \lor B)$ is retrieved
    \item[] Substituting $A := P$ and $B := P$ yields $P \rightarrow (P \lor P)$
\end{itemize}

\item \textbf{Complete the equivalence:}
LT completes the chain by linking the proved implications back to the original goal:
\begin{quote}
\begin{verbatim}
GIVEN       4.20      A=A
FORWARD CHAINING    4.25      P=(PVP)
Q.E.D.0
\end{verbatim}
\end{quote}
\textit{Logic form:} 
\begin{itemize}
    \item[] Because both $P \rightarrow (P \lor P)$ and $(P \lor P) \rightarrow P$ are established, the equivalence $P = (P \lor P)$ holds.
\end{itemize}

\end{enumerate}

\textbf{Performance Metrics:}
\begin{quote}
\begin{verbatim}
EFFORT              LIMIT 20000         ACTUAL 8727
SUBPROBLEMS         LIMIT 50            ACTUAL 2
SUBSTITUTIONS       LIMIT 50            ACTUAL 3
\end{verbatim}
\end{quote}

\medskip

\noindent The successful execution of complex derivations such as this
one, which use previously proved theorems, demonstrates the Logic
Theorist's capacity for a primitive form of machine learning: by
treating previously proved theorems, such as 2.20, as established
``givens'' for subsequent search paths, the system bootstraps its own
knowledge base, mirroring the cumulative nature of algebraic
reasoning.

\end{document}
